% ==================================================================
\section{模型评价}\label{sec:eval}
% ==================================================================

\subsection{模型的优点}

\begin{enumerate}[label=\arabic*.]
\item \textbf{三种对齐共用一个算子、一个接口。}问题一的三种方案只在“时间轴怎么定义”上不同，输出形状完全一致，下游可以直接切换；其中口径H复刻了附件2的词级对齐，每个位置都能回指到原文的词。
\item \textbf{有效时段的判据更贴合情感表达。}说话者人脸跟踪方案以“说话者的脸是否可见”划定有效时段，既剔除了无人时段，又保留了不说话时的表情；按声音判定只剩32\%的时长，而该判据平均保留93\%，语音、视觉非零位置分别为48.0和47.5。
\item \textbf{先查数据，再定方法。}建模前逐帧核查了附件2—附件4，发现训练数据与专项测试集的缺失形态根本不同，并据此改写了缺失模拟方式。仅这一项（M1）不改网络结构，就让段缺失 F1 提升了9.4个点。同一轮核查还确认了帧与词元的一一对应，为问题三的精确回指打下了基础。
\item \textbf{模块去留由数据决定。}我们为局部连续缺失准备了5个候选模块，按统一掩码逐级叠加筛选，最终只保留2个。被淘汰模块的结果和筛选过程中走过的弯路（在坏配置上叠加导致一轮结果作废）都完整保留在正文中。
\item \textbf{改进落在最难的情形上。}LDSD-Seg 在长段缺失上的 MAE 降幅（0.066）大于整体（0.047），验证集上完整模态与段缺失之间的 F1 落差由15.6个点缩小到0.1个点，鲁棒性的提升不是靠容易样本刷出来的。
\item \textbf{规律分析有自检。}三因子扫描前设置了“语言全缺必须最差”的校验项，靠它发现并修正了两处会导致错误结论的设计问题，避免了把 bug 当成发现写进报告。
\item \textbf{解释对得上、查得到。}解释层不改网络，融合层重写后与原层的最大输出差为 $4.77\times10^{-7}$；文本证据精确回指到词元，可当场人工核对；忠实性、稳定性两项指标表现良好，单样本归因与整体消融两条路线的结论一致。
\end{enumerate}

\subsection{模型的缺点}

\begin{enumerate}[label=\arabic*.]
\item \textbf{问题一的对齐依赖“词匀速发音”近似。}方法一、二在没有逐词时间戳的情况下按均匀发音给词元分配时刻，语速变化大时会有词级误差；方法三的说话者身份靠人脸面积匹配，面积比阈值0.4是经验值，没有逐样本自适应。
\item \textbf{对齐质量没有在下游任务上检验。}问题一的语音、视觉特征与附件2不在同一个特征空间，附件1又只有100条，我们只用非零位置数等结构指标评价了三种方案，没有比较它们对情感预测的实际影响；视觉特征也只用了人脸几何与像素统计量，对细微表情的刻画偏粗。
\item \textbf{语音、视觉证据只能定位到词级。}数据没有逐词时间戳，本文没有利用附件4的原始视频做语音—文本强制对齐，因此给不出秒级的语音时段和视觉关键帧。
\item \textbf{“重建有益”只得到部分复现。}在整模态缺失下，去掉重建或去掉 AHL 后 MAE 反而略好。本文给出了可能的解释，但由于时间有限，没有设计专门的实验去验证。
\item \textbf{结果多为单次运行。}同一个 M1+M4 配置在筛选实验与确认实验中的段缺失 F1 相差约1.5个点，说明存在运行间波动，而本文没有做多随机种子重复、也没有给出置信区间。
\item \textbf{训练轮数没有调优。}验证曲线显示模型在训练中后期出现过拟合，固定训练200轮并非最优，本文没有系统比较早停的效果。
\item \textbf{专项测试集缺少标签。}附件3、4的可信度只能借助附件2上的表现间接判断；强度被放大的现象只在附件4的5条可找回样本上观察到，样本太少。
\end{enumerate}

\subsection{改进方向}

针对上面的不足，后续可以从几个方面改进。第一，引入语音识别与强制对齐工具，得到逐词时间戳。这一步能同时改进两处：问题一的方法一、二不再依赖“词匀速发音”的近似；问题三的词级锚点可以换算成秒级的语音时段和视频关键帧，证据定位直接落到原始素材上。视觉方面，可以在人脸框内再接一个面部动作单元或表情特征提取器，补足对细微表情的刻画。第二，按验证集指标早停，并用多个随机种子重复主要实验，报告均值和区间，让小幅差异的比较更有依据。第三，在验证集上检查强度放大是否普遍存在，若存在则拟合线性校准。第四，模型对文本的依赖过重，一旦转写出错就容易被带偏，可以在训练时提高文本连续段缺失的比例，迫使模型更多地利用语音和视觉；中段缺失改善最小的原因也值得进一步分析。
